\subsection{Functoriality of the Poincaré groupoid}

Let X, Y be topological spaces and let \(f: X \to Y\) be a continuous map. The map \(c \mapsto f \circ c\) is a continuous map, denoted \(\Lambda(f)\) , from \(\Lambda(X) = \mathcal{C}_c(\mathbf{I}; X)\) to \(\Lambda(Y) = \mathcal{C}_c(\mathbf{I}; Y)\) (I, p.~132, lemma). It defines, by passing to subsets, the continuous maps

\[
\Lambda_ {x} (f) \colon \Lambda_ {x} (\mathrm{X}) \to \Lambda_ {f (x)} (\mathrm{Y}),
\]

continuous for \(x \in X\) ,

\[
\Lambda_ {x, y} (f) \colon \Lambda_ {x, y} (\mathrm{X}) \to \Lambda_ {f (x), f (y)} (\mathrm{Y}),
\]

continuous for \(x, y \in \mathrm{X}\) and

\[
\Omega (f) \colon \Omega (\mathrm{X}) \to \Omega (\mathrm{Y}).
\]

For \(x \in X\) , the map \(\Lambda_{x,x}(f)\) is also denoted \(\Omega_x(f)\) . By passing to arc-connected components (III, p.~290), we deduce from the map \(\Lambda_{x,y}(f)\) a map

\[
\varpi_ {x, y} (f) \colon \varpi_ {x, y} (\mathrm{X}) \to \varpi_ {f (x), f (y)} (\mathrm{Y}).
\]

Let \(x, y, z\) be points of X and let \(c \in \Lambda_{x,y}(X)\) and \(d \in \Lambda_{y,z}(X)\) be paths; by definition of the juxtaposition of paths, we have

\[
f \circ (c * d) = (f \circ c) * (f \circ d).
\]

By passing to strict homotopy classes, this relation follows.

\[
\varpi_ {x, z} (f) (\gamma \delta) = (\varpi_ {x, y} (f) (\gamma)) (\varpi_ {y, z} (f) (\delta))
\]

 for all \(\gamma \in \varpi_{x,y}(\mathrm{X})\) and every \(\delta \in \varpi_{y,z}(\mathrm{X})\) . Thus, the continuous map \(f\) and the maps \(\varpi_{x,y}(f)\) , for \(x\) and \(y \in \mathrm{X}\) , define a morphism of the groupoid \(\varpi(\mathrm{X})\) in the groupoid \(\varpi(\mathrm{Y})\) (II, p.~161, def. 3). It is denoted \(\varpi(f)\) and is called the Poincaré groupoid morphism deduced from the continuous map \(f\) . In particular, if \(x\) is a point of X, the map \(\varpi_{x,x}(f)\) is a homomorphism of the group \(\pi_1(\mathrm{X}, x)\) in the group \(\pi_1(\mathrm{Y}, f(x))\) ; this homomorphism is also denoted \(\pi_1(f, x)\) .

Remark 1. --- The homomorphism \(\varpi_{x,y}(f)\) is continuous if one equips the sets \(\varpi_{x,y}(\mathrm{X})\) and \(\varpi_{f(x),f(y)}(\mathrm{Y})\) with the quotient topology of the compact-open topology on \(\mathcal{C}_c(\mathbf{I};\mathrm{X})\) and \(\mathcal{C}_c(\mathbf{I};\mathrm{Y})\) .

To simplify the notation, if \(x\) , \(y \in X\) and \(\gamma \in \varpi_{x,y}(X)\) , we will sometimes write \(f_{*}(\gamma)\) for the element \(\varpi_{x,y}(f)(\gamma)\) of \(\varpi_{f(x),f(y)}(Y)\) .

Let X, Y, Z be topological spaces, let \(f: X \to Y\) and \(g: Y \to Z\) be continuous maps. For every path c in X, we have \((g \circ f) \circ c = g \circ (f \circ c)\) . It follows that one has

\[
\varpi (g \circ f) = \varpi (g) \circ \varpi (f).
\]

Let X and Y be topological spaces, let \(f_{0}\) and \(f_{1}\) be continuous maps from X to Y, and let \(\sigma\colon X\times I\to Y\) be a homotopy connecting \(f_{0}\) to \(f_{1}\) . The map \((c,t)\mapsto c(t)\) from \(\mathcal{C}_{c}(\mathbf{I};\mathrm{X})\times\mathbf{I}\) to X (III, p.~257, prop. 1) being continuous, the map from \(\mathcal{C}_{c}(\mathbf{I};\mathrm{X})\times\mathbf{I}\times\mathbf{I}\) to Y given by \((c,t,s)\mapsto\sigma(c(t),s)\) is also continuous. Consequently, the map \(\Sigma\colon(c,s)\mapsto\sigma(c(\cdot),s)\) is a continuous map from \(\mathcal{C}_{c}(\mathbf{I};\mathrm{X})\times\mathbf{I}\) to \(\mathcal{C}_{c}(\mathbf{I};\mathrm{Y})\) (loc. cit.). The map \(\Sigma\) is a homotopy connecting the maps \(\Lambda(f_{0})\) and \(\Lambda(f_{1})\) . By restriction to loop spaces, the map \(\Sigma\) induces a homotopy \(\Omega(\mathrm{X})\times\mathbf{I}\to\Omega(\mathrm{Y})\) connecting \(\Omega(f_{0})\) to \(\Omega(f_{1})\) . Let x be a point of X; suppose that \(\sigma\) is a pointed homotopy at x and set \(y=f_{0}(x)=f_{1}(x)\) . The map \(\Sigma\) then induces a continuous map from \(\Omega_{x}(\mathrm{X})\times\mathbf{I}\) to \(\Omega_{y}(\mathrm{Y})\) which is a pointed homotopy at \(e_{x}\) , connecting the maps \(\Omega_{x}(f_{0})\) and \(\Omega_{x}(f_{1})\) .

PROPOSITION 3. --- Let X and Y be topological spaces, \(f_{0}\) and \(f_{1}\) be continuous maps from X to Y, and let \(\sigma\colon\mathbf{X}\times\mathbf{I}\to\mathbf{Y}\) be a homotopy connecting \(f_{0}\) to \(f_{1}\) . Let x be a point of X; set \(y_{0}=f_{0}(x)\) , \(y_{1}=f_{1}(x)\) and denote by \(\delta\in\varpi_{y_{0},y_{1}}(\mathbf{Y})\) the class of the path d defined by \(d(t)=\sigma(x,t)\) for \(t\in\mathbf{I}\) . For all \(\gamma\in\pi_{1}(\mathbf{X},x)\) , we have \((f_{1})_{*}(\gamma)=\delta^{-1}\big((f_{0})_{*}(\gamma)\big)\delta\) .

Let c be a loop in X at x and let \(\gamma\) be its strict homotopy class. Set \(\gamma_{0} = (f_{0})_{*}(\gamma)\) and \(\gamma_{1} = (f_{1})_{*}(\gamma)\) ; these are the strict homotopy classes of \(f_{0} \circ c\) and \(f_{1} \circ c\) . For \((t, s) \in \mathbf{I} \times \mathbf{I}\) , set \(\varphi(t, s) = \sigma(c(t), s)\) . For all \(t \in I\) , we have \(\varphi(t, 0) = (f_{0} \circ c)(t)\) , \(\varphi(t, 1) = (f_{1} \circ c)(t)\) and \(\varphi(0, t) = \varphi(1, t) = d(t)\) . The relation \(\gamma_{0}\delta = \delta\gamma_{1}\) therefore follows from the following lemma.

Lemma 1. --- Let Y be a topological space and let \(\varphi\colon\mathbf{I}\times\mathbf{I}\to\mathbf{Y}\) be a continuous map. For \(t\in\mathbf{I}\) , set \(c_{0}(t)=\varphi(t,0)\) , \(c_{1}(t)=\varphi(t,1)\) , \(d_{0}(t)=\varphi(0,t)\) and \(d_{1}(t)=\varphi(1,t)\) . The paths \(c_{0}*d_{1}\) and \(d_{0}*c_{1}\) are strictly homotopic.

Let c denote the path in \(I \times I\) obtained by juxtaposing the paths \(t \mapsto (t,0)\) and \(t \mapsto (1,t)\) ; let d denote the path in \(I \times I\) obtained by juxtaposing the paths \(t \mapsto (0, t)\) and \(t \mapsto (t, 1)\) . The paths \(c\) and \(d\) have the same origin \((0, 0)\) and the same endpoint \((1, 1)\) . The map \((t, s) \mapsto (1 - s)c(t) + sd(t)\) from \(\mathbf{I} \times \mathbf{I}\) to \(\mathbf{I} \times \mathbf{I}\) is a strict homotopy joining \(c\) to \(d\) . We have \(c_0 * d_1 = \varphi \circ c\) and \(d_0 * c_1 = \varphi \circ d\) ; these two paths are therefore strictly homotopic.

COROLLARY 1. --- Let X and Y be topological spaces and let x be a point of X. Let \(f_0\) and \(f_1\) be continuous maps from X to Y. If there exists a pointed homotopy at x connecting \(f_0\) to \(f_1\) , we have \(\pi_1(f_0, x) = \pi_1(f_1, x)\) .

Let \(\sigma\) be a pointed homotopy at \(x\) connecting \(f_0\) to \(f_1\) . With the notations of Prop. 3, \(\delta\) is the class of a constant path with image \(f_0(x) = f_1(x)\) . The assertion follows from this.

COROLLARY 2. --- Let X and Y be topological spaces and let \(f\colon X \to Y\) be a homotopy equivalence. For every point \(x\) of X, the homomorphism

\[
\pi_ {1} (f, x) \colon \pi_ {1} (\mathrm{X}, x) \to \pi_ {1} (\mathrm{Y}, f (x))
\]

is an isomorphism.

Let \(g\) be a continuous map from Y to X, inverse up to homotopy of \(f\) . Let \(x\) be a point of X. It follows from prop. 3 applied to the homotopic maps \(\mathrm{Id}_{\mathrm{X}}\) and \(g \circ f: \mathrm{X} \to \mathrm{X}\) that the map \(\pi_1(g \circ f, x)\) is an isomorphism from the group \(\pi_1(\mathrm{X}, x)\) to the group \(\pi_1(\mathrm{X}, g \circ f(x))\) . Since

\[
\pi_ {1} (g \circ f, x) = \pi_ {1} (g, f (x)) \circ \pi_ {1} (f, x),
\]

the homomorphism \(\pi_{1}(f,x)\) is injective and the homomorphism \(\pi_{1}(g,f(x))\) is surjective. Since the map g is also a homotopy equivalence, the homomorphism \(\pi_{1}(g,f(x))\) is injective; it is therefore an isomorphism. Consequently, \(\pi_{1}(f,x)\) is an isomorphism.

Example. --- Let G be a topological group, and let e be its identity element. For every point g of G, the left and right translations, \(x \mapsto gx\) and \(x \mapsto xg\) are homeomorphisms of G onto itself (TG, III, p.~2) which map e to g. By Corollary 2, they induce isomorphisms of \(\pi_{1}(G, e)\) onto \(\pi_{1}(G, g)\) . These isomorphisms are not necessarily equal (IV, p.~459, exerc. 1).

COROLLARY 3. --- Let X be a topological space homotopy equivalent to a point. For every point x of X, the group \(\pi_{1}(X,x)\) is reduced to the identity element.

Cor. 3 applies in particular when X is the n-dimensional numerical space \(R^{n}\) and more generally when X is a subset of the space \(R^{n}\) which is star-shaped (III, p.~234) with respect to one of its points.

PROPOSITION 4. --- Let X be the product space of a family \((\mathrm{X}_{j})_{j\in\mathrm{J}}\) of topological spaces. The morphism of the groupoid \(\varpi(\mathrm{X})\) into the product of the groupoids \(\varpi(\mathrm{X}_{j})\) , for \(j\in J\) , defined by the family of morphisms \((\varpi(\mathrm{pr}_{j}))_{j\in\mathrm{J}}\) is an isomorphism.

Denote by \(\varphi\) this morphism of groupoids. The map deduced from it by passing to vertices is the identity map \(X \to \prod_{j} X_{j}\) . Let \(x = (x_{j})\) and \(y = (y_{j})\) be two points of \(X\) . Denote by \(\varphi_{x,y}\) the map from \(\varpi_{x,y}(X)\) to \(\prod_{j} \varpi_{x_{j},y_{j}}(X_{j})\) deduced from \(\varphi\) . If for every \(j \in J\) , \(c_{j}: I \to X_{j}\) is a path connecting \(x_{j}\) to \(y_{j}\) , the map \(t \mapsto (c_{j}(t))\) is a path in \(X\) connecting \(x\) to \(y\) (TG, I, p.~25, prop. 1). This proves that \(\varphi_{x,y}\) is surjective. Let \(c\) and \(d\) be two paths in \(X\) connecting \(x\) to \(y\) . Suppose that for every \(j \in J\) there exists a strict homotopy \(\sigma_{j}: I \times I \to X_{j}\) connecting \(\operatorname{pr}_{j} \circ c\) to \(\operatorname{pr}_{j} \circ d\) . Then, the map \((t,s) \mapsto (\sigma_{j}(t,s))\) from \(I \times I\) to \(X\) is a strict homotopy joining \(c\) to \(d\) (loc. cit.). This proves that \(\varphi_{x,y}\) is injective.

COROLLARY. --- Let \(x = (x_{j})_{j \in \mathrm{J}}\) be a point of X. The map

\[
\pi_ {1} (X, x) \to \prod_ {j \in J} \pi_ {1} (X _ {j}, x _ {j})
\]

deduced from the maps \(\pi_1(\mathrm{pr}_j, x_j)\) is a group isomorphism.

This isomorphism is said to be canonical. In what follows, we will often identify \(\pi_{1}(\mathrm{X},x)\) with \(\prod_{j\in\mathrm{J}}\pi_{1}(\mathrm{X}_{j},x_{j})\) by means of this isomorphism.

Remarks. --- 2) Let \((\mathrm{X}_{j})_{j\in\mathrm{J}}\) be a family of topological spaces. Let X denote the product topological space \(\prod_{j\in J}X_{j}\) and let \(x=(x_{j})\) be a point of X.

For every \(i \in \mathrm{J}\) , let \(u_i: \mathrm{X}_i \to \mathrm{X}\) be the map such that, for \(z \in \mathrm{X}_i\) , \(\operatorname{pr}_i \circ u_i(z) = z\) and, for all \(j \in \mathrm{J}\) distinct from \(i\) , \(\operatorname{pr}_j \circ u_i(z) = x_j\) . The map \(u_i\) is continuous and the map \(\pi_1(u_i, x_i)\) is identified with the canonical injection of the factor \(\pi_1(\mathrm{X}_i, x_i)\) in the product group of the family \((\pi_1(\mathrm{X}_j, x_j))_{j \in \mathrm{J}}\) (cf.~A, I, p.~45).

Suppose that the set J is finite and, for every \(j \in J\) , let \(\gamma_{j}\) be an element of \(\pi_{1}(X_{j}, x_{j})\) . The element \((\gamma_{j})\) of \(\pi_{1}(X, x)\) is the composite of the loop classes \((u_{j})_{*}(\gamma_{j})\) , \(j \in J\) , which are pairwise permutable.

\begin{enumerate}
\def\labelenumi{\arabic{enumi})}
\setcounter{enumi}{2}
\tightlist
\item
  Let \((\mathrm{X}_j)_{j\in \mathrm{J}}\) be a family of topological spaces. Let X denote the product topological space \(\prod_{j\in \mathrm{J}}\mathrm{X}_j\) and let \(x = (x_{j})\) be a point of X.
\end{enumerate}

Endow the sets \(\pi_1(\mathrm{X},x)\) and \(\pi_1(\mathrm{X}_j,x_j)\) with the quotient topology of the compact-open topology on the spaces \(\Lambda_x(\mathrm{X})\) and \(\Lambda_{x_j}(\mathrm{X}_j)\) . The isomorphism \(\pi_1(\mathrm{X},x)\to \prod_{j\in \mathrm{J}}\pi_1(\mathrm{X}_j,x_j)\) is then a homeomorphism. It is continuous (III, p.~294, remark 1). The compact-open topology on \(\Lambda (\mathrm{X})\) is generated by the subsets of the form \(T(\mathrm{K},\mathrm{U})\) , where K is a compact subset of I and U is an open subset of X. For \(j\in \mathrm{J}\) , let \(\mathrm{U}_j\) be an open subset of \(\mathrm{X}_j\) such that \(\prod_{j\in \mathrm{J}}\mathrm{U}_j\subset \mathrm{U}\) . Then \((\mathrm{pr}_j)_*(T(\mathrm{K},\mathrm{U}))\) contains \(T(\mathrm{K},\mathrm{U}_j)\) . This shows that the maps \((\mathrm{pr}_j)_*: \Lambda_x(\mathrm{X})\to \Lambda_{x_j}(\mathrm{X}_j)\) are open, and the maps \(\pi_1(\mathrm{pr}_j,x_j)\) are also open. Since they are surjective, the map \(\pi_1(\mathrm{X},x)\to \prod_{j\in \mathrm{J}}\pi_1(\mathrm{X}_j,x_j)\) is open (TG, I, p.~34, prop. 8). Being continuous and bijective, it is a homeomorphism (TG, I, p.~30, example 2).

PROPOSITION 5. — Let X be a topological space and let \((\mathrm{A}_{i})_{i\in\mathrm{I}}\) be an increasing family of subsets of X, indexed by a filtered ordered set I, such that every quasi-compact subset of X is contained in one of the \(A_{i}\) . The canonical groupoid morphism

\[
\rho \colon \varinjlim_ {i \in \mathrm{I}} \varpi (\mathrm{A} _ {i}) \to \varpi (\mathrm{X}),
\]

deduced from the canonical injections of \(A_{i}\) into X, is an isomorphism. If \(i \leqslant j\) , denote by \(\rho_{j,i}\) the groupoid morphism \(\varpi(\mathrm{A}_{i}) \to \varpi(\mathrm{A}_{j})\) deduced from the injection of \(A_{i}\) into \(A_{j}\) . Since the map deduced from \(\rho_{j,i}\) by passing to vertices is the injection \(A_{i} \to A_{j}\) and the \(A_{i}\) cover X, the map induced from \(\rho\) by passing to vertices is bijective.

Let \(a\) and \(b\) be points of X and let \(c\) be a path connecting \(a\) to \(b\) in X. The image of \(c\) is a quasi-compact subset of X (TG, I, p.~62, th. 2), since \(\mathbf{I}\) is compact. There therefore exists an element \(i \in \mathrm{I}\) such that the image of \(c\) is contained in \(\mathrm{A}_i\) . Consequently, the map deduced from \(\rho\) by passing to arrow sets is surjective.

Let \(i \in \mathrm{I}\) , let \(a\) and \(b\) be points of \(\mathrm{A}_i\) , let \(c\) and \(c'\) be paths joining \(a\) to \(b\) in \(\mathrm{A}_i\) ; let \(h\) be a strict homotopy joining \(c\) to \(c'\) in \(\mathrm{X}\) . Since \(\mathbf{I} \times \mathbf{I}\) is compact, \(h(\mathbf{I} \times \mathbf{I})\) is a quasi-compact subset of \(\mathrm{X}\) (loc. cit.) and there exists an element \(i \in \mathrm{I}\) such that the image of \(h\) is contained in \(\mathrm{A}_i\) . The paths \(c\) and \(c'\) are strictly homotopic in \(A_{i}\) ; a fortiori, the path classes {[}c{]} and \([c']\) have the same image in \(\varinjlim\varpi(A_{i})\) . Consequently, \(\rho\) is injective.

COROLLARY. --- Let a be a point of X and let J be the set of \(i \in I\) such that \(a \in A_{i}\) . The canonical homomorphism

\[
\varinjlim_ {i \in J} \pi_ {1} (A _ {i}, a) \to \pi_ {1} (X, a)
\]

is bijective.

Remark 4. --- The proposition and its corollary apply in particular when \((\mathrm{A}_{i})_{i\in\mathrm{I}}\) is an increasing family of subsets of X indexed by a filtered ordered set I such that the interiors of the \(A_{i}\) cover X.

